\chapter{方法}\label{chap:method}

本章介绍基于物理信息神经网络（PINN）与大语言模型（LLM）的统一期权定价框架。系统整体流程为：用户以自然语言描述期权需求，LLM 路由层解析参数并选择模型，对应的 PINN 求解器返回定价结果。

% ─────────────────────────────────────────────────────────────────────────────
\section{PINN 统一框架}\label{sec:pinn_framework}

\subsection{问题形式化}

设期权价值 $V$ 满足如下一般形式的 PDE：
\begin{equation}\label{eq:general_pde}
  \mathcal{F}[V](x, t) = 0, \quad (x, t) \in \Omega \times [0, T),
\end{equation}
其中 $x$ 为状态变量（BSM/CEV 中为股价 $S$，Heston 中为 $(S, v)$），$\mathcal{F}$ 为微分算子，$T$ 为到期日。边界条件与终值条件分别为：
\begin{align}
  V(x, T) &= \phi(x), \quad x \in \Omega, \label{eq:terminal} \\
  \mathcal{B}[V](x, t) &= g(x, t), \quad x \in \partial\Omega. \label{eq:boundary}
\end{align}
对于欧式看涨期权，终值条件为 $\phi(x) = \max(S - K, 0)$。

\subsection{损失函数}

PINN 用神经网络 $\hat{V}_\theta(x, t)$ 近似真实解 $V(x, t)$，通过最小化以下复合损失函数训练网络参数 $\theta$：
\begin{equation}\label{eq:loss}
  \mathcal{L}(\theta) = w_\text{pde}\,\mathcal{L}_\text{pde}
                      + w_\text{ic}\,\mathcal{L}_\text{ic}
                      + w_\text{bc}\,\mathcal{L}_\text{bc},
\end{equation}
其中各项定义为：
\begin{align}
  \mathcal{L}_\text{pde} &= \frac{1}{N_c}\sum_{i=1}^{N_c}
    \left|\mathcal{F}[\hat{V}_\theta](x_i, t_i)\right|^2, \\
  \mathcal{L}_\text{ic}  &= \frac{1}{N_T}\sum_{j=1}^{N_T}
    \left|\hat{V}_\theta(x_j, T) - \phi(x_j)\right|^2, \\
  \mathcal{L}_\text{bc}  &= \frac{1}{N_b}\sum_{k=1}^{N_b}
    \left|\mathcal{B}[\hat{V}_\theta](x_k, t_k) - g(x_k, t_k)\right|^2.
\end{align}
配点 $\{(x_i, t_i)\}_{i=1}^{N_c}$ 在求解域内随机采样，无需网格剖分。PDE 残差通过自动微分（automatic differentiation）精确计算。

\subsection{硬约束终值条件（ICPINN）}

软约束方式将终值条件纳入损失函数，网络只能近似满足。参考 Guo 等人\citep{guo2024icpinn}的 ICPINN 方法，对 Heston 模型采用输出变换将终值条件硬编码：
\begin{equation}\label{eq:hard_constraint}
  \hat{V}_\theta(S, v, t) = \phi(S) + (T - t)\cdot\text{net}_\theta(S, v, t).
\end{equation}
当 $t = T$ 时，$(T-t) = 0$，故 $\hat{V}_\theta(S, v, T) = \phi(S)$ 精确成立，无需在损失函数中加入 $\mathcal{L}_\text{ic}$ 项。

% ─────────────────────────────────────────────────────────────────────────────
\section{网络架构}\label{sec:architecture}

\subsection{门控网络（BSM 与 CEV）}

BSM 和 CEV 模型采用门控网络结构\citep{dhiman2023pinn}，由浅层分支与深层分支并联组成：
\begin{align}
  h_s &= f_\text{shallow}(x), \\
  h_d &= f_\text{deep}(x), \\
  g   &= \sigma(W_g h_s + b_g), \\
  h   &= g \odot h_s + (1-g) \odot h_d, \\
  \hat{V} &= W_\text{out} h + b_\text{out},
\end{align}
其中 $\sigma(\cdot)$ 为 sigmoid 函数，$\odot$ 为逐元素乘法，$g \in (0,1)$ 为门控权重。浅层分支（1 个隐藏层，64 个神经元）捕捉期权价格的线性部分（深度实值/虚值区域），深层分支（4 个隐藏层，64 个神经元）捕捉平值附近的非线性特征。激活函数均为 $\tanh$。

\subsection{深层全连接网络（Heston）}

Heston 模型输入维度为 3（$S, v, t$），采用对数价格变换与 Fourier 特征嵌入的改进架构，并结合式~\eqref{eq:hard_constraint}~的硬约束输出变换。

\textbf{对数价格变换}：将股价输入从 $S \in [0, S_\max]$ 变换为对数坐标 $x = \ln(S/K)$，使 PDE 扩散项从 $\frac{1}{2}vS^2 V_{SS}$ 变为 $\frac{1}{2}v V_{xx}$，消除了 $S^2$ 量级失衡。在对数坐标下，Heston PDE 变为：
\begin{equation}\label{eq:heston_logprice}
  \frac{\partial V}{\partial t}
  + \frac{1}{2}v\frac{\partial^2 V}{\partial x^2}
  + \left(r - \frac{1}{2}v\right)\frac{\partial V}{\partial x}
  + \rho\xi v\frac{\partial^2 V}{\partial x\partial v}
  + \frac{1}{2}\xi^2 v\frac{\partial^2 V}{\partial v^2}
  + \kappa(\theta-v)\frac{\partial V}{\partial v}
  - rV = 0.
\end{equation}

\textbf{Fourier 特征嵌入}：在输入层加入随机 Fourier 特征映射 $\gamma(x) = [\sin(Bx), \cos(Bx)]$（$B \in \mathbb{R}^{3 \times n_f}$ 为随机矩阵，$n_f=16$），提升网络对高频分量的表达能力，有助于捕捉平值附近的非线性特征。

\textbf{网络结构}：8 层全连接网络，每层 128 个神经元，激活函数为 $\tanh$，输入为 Fourier 嵌入后的 32 维特征向量。

\textbf{自适应损失权重}：训练过程中每 1000 步根据各损失项的梯度范数动态调整权重（Wang 等人\citep{wang2022ntk}），防止边界条件损失被 PDE 残差损失淹没：
\begin{equation}
  w_i \leftarrow \frac{\sum_j \|\nabla_\theta \mathcal{L}_j\|}{\|\nabla_\theta \mathcal{L}_i\|}.
\end{equation}

\textbf{RAR 自适应采样}：每 2000 步根据当前 PDE 残差大小，在高残差区域（深度虚值 $S \ll K$）补充配点，强制网络关注接近零的期权价格区域（Lu 等人\citep{lu2021deepxde}）。

\subsection{输入归一化}

为提升训练稳定性，所有输入在送入网络前进行归一化：
\begin{equation}
  \tilde{S} = \frac{S}{S_\max}, \quad
  \tilde{v} = \frac{v}{v_\max}, \quad
  \tilde{t} = \frac{t}{T},
\end{equation}
使各输入均落在 $[0, 1]$ 区间，避免不同量级的输入导致梯度不平衡。

% ─────────────────────────────────────────────────────────────────────────────
\section{波动率、隐含波动率与希腊字母}\label{sec:greeks}

\subsection{历史波动率与隐含波动率}

\textbf{历史波动率}（Historical Volatility，HV）是标的资产对数收益率的样本标准差，反映过去一段时间内价格的实际波动幅度：
\begin{equation}
  \sigma_\text{HV} = \sqrt{\frac{252}{n-1}\sum_{i=1}^{n}\left(\ln\frac{S_i}{S_{i-1}} - \bar{\mu}\right)^2},
\end{equation}
其中 252 为年化因子（交易日数），$\bar{\mu}$ 为样本均值。历史波动率是向后看的（backward-looking），只能描述已发生的价格变动。

\textbf{隐含波动率}（Implied Volatility，IV）是使 BSM 模型价格等于市场观测价格时，反推得到的波动率参数：
\begin{equation}\label{eq:iv}
  \sigma_\text{IV} = \text{BSM}^{-1}(V^\text{mkt};\, S, K, T, r),
\end{equation}
即求解方程 $V_\text{BS}(S, K, T, r, \sigma) = V^\text{mkt}$ 中的 $\sigma$。隐含波动率是向前看的（forward-looking），反映市场对未来波动率的预期，是期权市场的核心报价语言。

实践中，同一标的不同行权价和到期日的期权往往对应不同的隐含波动率，形成\textbf{波动率微笑}（Volatility Smile）或\textbf{波动率偏斜}（Volatility Skew）。BSM 假设常数波动率，无法解释这一现象；CEV 和 Heston 模型正是为此而设计。

\subsection{希腊字母}

期权的\textbf{希腊字母}（Greeks）描述期权价格对各市场参数的敏感度，是风险管理和对冲的核心工具。主要希腊字母定义如下：

\begin{align}
  \Delta &= \frac{\partial V}{\partial S}, \label{eq:delta} \\
  \Gamma &= \frac{\partial^2 V}{\partial S^2}, \label{eq:gamma} \\
  \Theta &= \frac{\partial V}{\partial t}, \label{eq:theta} \\
  \mathcal{V} &= \frac{\partial V}{\partial \sigma}, \label{eq:vega} \\
  \rho_r &= \frac{\partial V}{\partial r}. \label{eq:rho_r}
\end{align}

各希腊字母的经济含义与 BSM 解析表达式汇总于表~\ref{tab:greeks}。

\begin{table}[htbp]
  \centering
  \caption{主要希腊字母的定义、含义与 BSM 解析式（欧式看涨期权）}
  \label{tab:greeks}
  \begin{tabular}{clll}
    \hline
    符号 & 定义 & 经济含义 & BSM 解析式 \\
    \hline
    $\Delta$ & $\partial V/\partial S$ & 股价变动 \$1 时期权价格变动量 & $\Phi(d_1)$ \\
    $\Gamma$ & $\partial^2 V/\partial S^2$ & Delta 对股价的变化率 & $\phi(d_1)/(S\sigma\sqrt{T})$ \\
    $\Theta$ & $\partial V/\partial t$ & 时间流逝对期权价值的侵蚀 & $-S\phi(d_1)\sigma/(2\sqrt{T}) - rKe^{-rT}\Phi(d_2)$ \\
    $\mathcal{V}$ & $\partial V/\partial \sigma$ & 波动率变动 1\% 时价格变动量 & $S\phi(d_1)\sqrt{T}$ \\
    $\rho_r$ & $\partial V/\partial r$ & 利率变动 1\% 时价格变动量 & $KTe^{-rT}\Phi(d_2)$ \\
    \hline
  \end{tabular}
\end{table}

其中 $\phi(\cdot)$ 为标准正态密度函数，$\Phi(\cdot)$ 为标准正态分布函数，$d_1, d_2$ 的定义见式~\eqref{eq:bs_d1}。

\subsection{Heston 模型下希腊字母的高效计算}

BSM 模型的希腊字母均有闭合解析式，计算代价极低。对于 CEV 模型，一般需要数值微分（有限差分）。Heston 模型的优势在于：其特征函数方法允许对 $S$、$v_0$、$r$ 等参数\textbf{解析求导}，从而一次积分同时得到多个希腊字母，无需重复调用定价函数。

具体地，Heston 模型的 Delta 和 Vega（对 $v_0$ 的敏感度）可直接从特征函数积分中读出：
\begin{align}
  \Delta_\text{Heston} &= \frac{\partial V}{\partial S} = P_1 + S\frac{\partial P_1}{\partial S} - Ke^{-rT}\frac{\partial P_2}{\partial S}, \label{eq:heston_delta} \\
  \mathcal{V}_\text{Heston} &= \frac{\partial V}{\partial v_0} = S\frac{\partial P_1}{\partial v_0} - Ke^{-rT}\frac{\partial P_2}{\partial v_0}, \label{eq:heston_vega}
\end{align}
其中 $\partial P_j/\partial S$ 和 $\partial P_j/\partial v_0$ 通过对特征函数的被积函数直接求导得到，计算量与定价本身相当\citep{heston1993closed}。这一特性使 Heston 模型在需要实时计算大量希腊字母的做市场景中具有显著优势。

在 PINN 框架下，希腊字母可通过自动微分（automatic differentiation）直接从网络输出对输入求导得到，无需额外推导解析式：
\begin{equation}
  \Delta_\text{PINN} = \frac{\partial \hat{V}_\theta}{\partial S}, \quad
  \Gamma_\text{PINN} = \frac{\partial^2 \hat{V}_\theta}{\partial S^2},
\end{equation}
这是 PINN 相对于传统有限差分方法的另一优势：希腊字母与期权价格在同一次前向传播中同时计算，且精度由自动微分保证，无有限差分的截断误差。

% ─────────────────────────────────────────────────────────────────────────────
\section{BSM-PINN}\label{sec:bsm_pinn}

BSM 模型的 PDE 算子为：
\begin{equation}
  \mathcal{F}_\text{BSM}[V] =
    \frac{\partial V}{\partial t}
    + \frac{1}{2}\sigma^2 S^2 \frac{\partial^2 V}{\partial S^2}
    + rS\frac{\partial V}{\partial S} - rV.
\end{equation}
边界条件设置为：
\begin{align}
  V(0, t) &= 0, \\
  V(S_\max, t) &= S_\max - Ke^{-r(T-t)}.
\end{align}
上边界采用深度实值近似：当 $S \gg K$ 时，期权几乎必然被执行，价值近似为远期价值 $S - Ke^{-r(T-t)}$。

BSM 模型存在解析解，可用于精确验证 PINN 的精度：
\begin{equation}
  V_\text{BS}(S, t) = S\Phi(d_1) - Ke^{-r(T-t)}\Phi(d_2),
\end{equation}
其中
\begin{equation}\label{eq:bs_d1}
  d_1 = \frac{\ln(S/K) + (r + \sigma^2/2)(T-t)}{\sigma\sqrt{T-t}}, \quad d_2 = d_1 - \sigma\sqrt{T-t},
\end{equation}
$\Phi(\cdot)$ 为标准正态分布函数。

BSM 模型的适用场景为：短期、流动性好、平值附近的期权，或需要快速计算隐含波动率和希腊字母的场景。其核心假设——常数波动率——在实践中往往不成立，但 BSM 解析解的简洁性使其成为期权市场的报价基准，所有其他模型的定价结果通常也以"BSM 隐含波动率"的形式表达，便于横向比较。

% ─────────────────────────────────────────────────────────────────────────────
\section{CEV-PINN}\label{sec:cev_pinn}

CEV 模型将 BSM 中的常数波动率替换为依赖股价的局部波动率，PDE 算子为：
\begin{equation}
  \mathcal{F}_\text{CEV}[V] =
    \frac{\partial V}{\partial t}
    + \frac{1}{2}\sigma^2 S^{2\beta} \frac{\partial^2 V}{\partial S^2}
    + rS\frac{\partial V}{\partial S} - rV.
\end{equation}
本文取 $\beta = 0.5$（平方根过程），此时波动率 $\sigma S^{-0.5}$ 随股价上涨而下降，能够捕捉股票市场中的杠杆效应。

CEV 模型一般无闭合解析解，本文采用 Crank-Nicolson 有限差分法生成参考价格用于验证。

CEV 模型的适用场景为个股期权，尤其是财务杠杆较高或股价较低的公司。其核心优势在于以最少的额外参数（仅增加 $\beta$）捕捉了杠杆效应，计算代价远低于 Heston 模型。本文固定 $\beta=0.5$，依据 Beckers\citep{beckers1980constant} 的实证共识和 Schroder\citep{schroder1989computing} 的解析解存在性结论。CEV 模型的局限在于波动率仍为确定性函数，无法描述波动率的随机聚集现象，因此不适用于指数期权的精确定价。

% ─────────────────────────────────────────────────────────────────────────────
\section{Heston-PINN}\label{sec:heston_pinn}

Heston 模型的 PDE 算子为：
\begin{equation}
  \mathcal{F}_\text{Heston}[V] =
    \frac{\partial V}{\partial t}
    + \frac{1}{2}vS^2\frac{\partial^2 V}{\partial S^2}
    + \rho\xi vS\frac{\partial^2 V}{\partial S\partial v}
    + \frac{1}{2}\xi^2 v\frac{\partial^2 V}{\partial v^2}
    + rS\frac{\partial V}{\partial S}
    + \kappa(\theta-v)\frac{\partial V}{\partial v}
    - rV.
\end{equation}
边界条件设置为：
\begin{align}
  V(0, v, t) &= 0, \\
  V(S_\max, v, t) &= S_\max - Ke^{-r(T-t)}, \\
  V(S, v_\max, t) &\approx S.
\end{align}
当 $v \to v_\max$ 时，极高波动率下期权价值趋近于标的资产价格本身。

Heston 模型具有半解析解（特征函数方法），本文用其作为验证基准：
\begin{equation}
  V_\text{Heston}(S, v, t) = S\,P_1 - Ke^{-r(T-t)}\,P_2,
\end{equation}
其中 $P_1, P_2$ 通过对特征函数的数值积分求得\citep{heston1993closed}。

Heston 模型是指数期权定价的行业标准，适用于需要精确拟合波动率微笑和波动率期限结构的场景。其五个参数均有明确的经济含义：$\kappa$ 控制波动率均值回归速度（$\kappa$ 越大，波动率越快回归长期均值 $\theta$）；$\xi$ 为波动率的波动率，决定波动率微笑的曲率；$\rho$ 捕捉股价与波动率的负相关性（杠杆效应），通常为负值。Heston 模型的主要挑战在于参数校准：五个参数需要同时从期权链数据中反推，优化问题存在多个局部极小值，对初始值敏感。本文采用 L-BFGS-B 算法进行校准，以市场期权价格的均方误差为目标函数。

% ─────────────────────────────────────────────────────────────────────────────
\section{LLM 路由层}\label{sec:llm_router}

\subsection{设计动机}

三种 PINN 模型分别适用于不同的市场假设：BSM 适用于波动率平稳的市场，CEV 适用于存在杠杆效应的股票市场，Heston 适用于波动率本身随机变化的场景。传统工具要求用户手动指定模型类型和参数，门槛较高。

本文引入 LLM 作为智能路由层，使用户可以通过自然语言描述期权需求，由 LLM 自动完成模型选择和参数提取。

\subsection{路由规则}

LLM 根据以下规则选择模型：
\begin{itemize}
  \item 若用户仅提供常数波动率 $\sigma$，选择 \textbf{BSM}；
  \item 若用户提及弹性参数 $\beta$ 或"CEV"，选择 \textbf{CEV}；
  \item 若用户提供随机波动率参数（$\kappa, \theta, \xi, \rho, v_0$）或提及"Heston"，选择 \textbf{Heston}；
  \item 默认选择 \textbf{BSM}。
\end{itemize}

\subsection{参数提取}

LLM 接收用户的自然语言输入，输出结构化 JSON 参数：
\begin{verbatim}
{
  "model": "BSM",
  "option_type": "call",
  "S": 100.0,
  "K": 100.0,
  "T": 1.0,
  "r": 0.05,
  "sigma": 0.2
}
\end{verbatim}
本文使用 DeepSeek API（兼容 OpenAI 接口），通过 `response_format: json_object` 约束输出格式，确保参数提取的可靠性。

\subsection{系统流程}

完整的定价流程如下：
\begin{enumerate}
  \item 用户输入自然语言描述（中文或英文均可）；
  \item LLM 解析输入，提取期权参数，选择定价模型；
  \item 对应的 PINN 模型加载预训练权重，对给定参数即时推断；
  \item 系统返回期权价格及所用模型信息。
\end{enumerate}

% ─────────────────────────────────────────────────────────────────────────────
\section{本章小结}\label{sec:method_summary}

本章介绍了基于 PINN 与 LLM 的统一期权定价框架。在 PINN 部分，三大模型共用统一的损失函数结构，Heston 模型额外采用 ICPINN 硬约束确保终值条件精确满足；在网络架构上，BSM/CEV 使用门控网络，Heston 使用深层全连接网络。在 LLM 部分，路由层通过结构化输出实现可靠的参数提取和模型选择。下一章将报告三大模型的实验结果。
